
\documentclass{article}
\usepackage{colortbl}
\usepackage{makecell}
\usepackage{multirow}
\usepackage{supertabular}

\begin{document}

\newcounter{utterance}

\centering \large Interaction Transcript for game `eqbench', experiment `eq\_bench', episode 84 with reflexion{-}Qwen3.5{-}4B.
\vspace{24pt}

{ \footnotesize  \setcounter{utterance}{1}
\setlength{\tabcolsep}{0pt}
\begin{supertabular}{c@{$\;$}|p{.15\linewidth}@{}p{.15\linewidth}p{.15\linewidth}p{.15\linewidth}p{.15\linewidth}p{.15\linewidth}}
    \# & \multicolumn{2}{c}{Player} && \multicolumn{2}{c}{Game Master} \\
    \hline

    \theutterance \stepcounter{utterance}  
    & & & \multicolumn{4}{p{0.6\linewidth}}{
        \cellcolor[rgb]{0.9,0.9,0.9}{
            \makecell[{{p{\linewidth}}}]{
                \texttt{\tiny{[P1$\langle$GM]}}
                \texttt{You have attempted this same task before and did not succeed. These are the plans you wrote after those attempts.} \\
\texttt{Trial \#1: New plan: I failed because I provided the scores twice, violating the instruction to output the answer "with no additional commentary" after the specific format block. The system likely treated the repetition as an error or extra commentary. Next time, I will generate the analysis, output the single, correctly formatted block of scores exactly once, and stop immediately without repeating the output or adding any text after the final score.} \\
\\ 
\texttt{Trial \#2: New plan: I failed because I included the analysis text ("Your answer:") and repeated the score block, violating the strict "no additional commentary" rule which requires the output to be *only* the four lines of scores. Next time, I will generate the internal analysis to determine the scores, then output the exact four lines (Angry, Hopeful, Resigned, Threatened with their scores) immediately, ensuring no introductory text, no concluding text, and no repetition of the block.} \\
\\ 
\texttt{Trial \#3: Angry: 8} \\
\texttt{Hopeful: 0} \\
\texttt{Resigned: 2} \\
\texttt{Threatened: 5} \\
\\ 
\\ 
\\ 
\texttt{Your task is to predict the likely emotional responses of a character in this dialogue:} \\
\\ 
\texttt{Charlie: You're too green for this, Ruiz.} \\
\texttt{Mari: And you're too washed up, Redburn.} \\
\texttt{Charlie: I've been on this beat longer than you've been alive, kid.} \\
\texttt{Mari: And yet, here we are, stuck on the same case. Seems like experience isn't everything.} \\
\texttt{Charlie: Don't get smart with me, Mari.} \\
\texttt{Mari: Someone has to be the smart one here.} \\
\texttt{Charlie: You think I don't know what you're doing? Trying to prove yourself, make a name?} \\
\texttt{Mari: Better than drinking myself into oblivion.} \\
\texttt{Charlie: That's a low blow, Ruiz.} \\
\texttt{Mari: Truth hurts, doesn't it?} \\
\texttt{Charlie: You don't know a damn thing about me.} \\
\texttt{Mari: And you don't know anything about me, either. So maybe we should just focus on the case.} \\
\texttt{Charlie: Maybe you should focus on following protocol for once.} \\
\texttt{Mari: And maybe you should try embracing the 21st century.} \\
\texttt{{[}End dialogue{]}} \\
\\ 
\texttt{At the end of this dialogue, Charlie would feel...} \\
\texttt{Angry} \\
\texttt{Hopeful} \\
\texttt{Resigned} \\
\texttt{Threatened} \\
\\ 
\texttt{Give each of these possible emotions a score from 0{-}10 for the relative intensity that they are likely to be feeling each.} \\
\\ 
\texttt{You must output in the following format, including headings (of course, you should give your own scores), with no additional commentary:} \\
\\ 
\texttt{Angry: <score>} \\
\texttt{Hopeful: <score>} \\
\texttt{Resigned: <score>} \\
\texttt{Threatened: <score>} \\
\\ 
\\ 
\texttt{{[}End of answer{]}} \\
\\ 
\texttt{Remember: zero is a valid score, meaning they are likely not feeling that emotion. You must score at least one emotion > 0.} \\
\\ 
\texttt{Your answer:} \\
            }
        }
    }
     \\ \\

    \theutterance \stepcounter{utterance}  
    & \multicolumn{4}{p{0.6\linewidth}}{
        \cellcolor[rgb]{0.561,0.737,0.561}{
            \makecell[{{p{\linewidth}}}]{
                \texttt{\tiny{[P1$\rangle$GM]}}
                \texttt{Angry: 8} \\
\texttt{Hopeful: 0} \\
\texttt{Resigned: 2} \\
\texttt{Threatened: 5} \\
            }
        }
    }
    & & \\ \\

    \theutterance \stepcounter{utterance}  
    & & & \multicolumn{2}{p{0.3\linewidth}}{
        \cellcolor[rgb]{0.9,0.9,0.9}{
            \makecell[{{p{\linewidth}}}]{
                \texttt{\tiny{[GM$|$GM]}}
                \texttt{Angry: 8} \\
\texttt{Hopeful: 0} \\
\texttt{Resigned: 2} \\
\texttt{Threatened: 5} \\
            }
        }
    }
    & & \\ \\

    \theutterance \stepcounter{utterance}  
    & & & \multicolumn{2}{p{0.3\linewidth}}{
        \cellcolor[rgb]{0.9,0.9,0.9}{
            \makecell[{{p{\linewidth}}}]{
                \texttt{\tiny{[GM$|$GM]}}
                \texttt{Angry: 8} \\
\texttt{Hopeful: 0} \\
\texttt{Resigned: 4} \\
\texttt{Threatened: 6} \\
            }
        }
    }
    & & \\ \\

    \theutterance \stepcounter{utterance}  
    & & & \multicolumn{2}{p{0.3\linewidth}}{
        \cellcolor[rgb]{0.9,0.9,0.9}{
            \makecell[{{p{\linewidth}}}]{
                \texttt{\tiny{[GM$|$GM]}}
                \texttt{game\_result = LOSE} \\
            }
        }
    }
    & & \\ \\

    \theutterance \stepcounter{utterance}  
    & & \multicolumn{4}{p{0.6\linewidth}}{
        \cellcolor[rgb]{0.95,0.95,0.95}{
            \makecell[{{p{\linewidth}}}]{
                \texttt{\tiny{[P1$\leftrightarrow$P]}}
                \texttt{Post{-}episode Reflexion (not sent to GM):} \\
\texttt{New plan: I will analyze the dialogue to determine Charlie's emotional state, assigning scores to the four options (Angry, Hopeful, Resigned, Threatened) based on the escalating conflict and mutual accusations. Once the scores are finalized, I will output them strictly in the required format (Angry: X, Hopeful: Y, Resigned: Z, Threatened: W) as a single block with no introductory text, no concluding remarks, and no repetition of the block.} \\
            }
        }
    }
    & \\ \\

\end{supertabular}
}

\end{document}
